% Luc Destombe --- licence CC BY-NC-SA 4.0
% https://creativecommons.org/licenses/by-nc-sa/4.0/deed.fr
% Fichier autonome : compiler avec pdflatex, deux fois (signets, nombre de pages).
\documentclass[a4paper,10pt]{article}
\makeatletter
\newif\iflycee@niveaux
\newif\iflycee@palatino
\lycee@palatinotrue

\RequirePackage[utf8]{inputenc}
\RequirePackage[T1]{fontenc}
\RequirePackage[french]{babel}
\RequirePackage{etoolbox}

\newcommand{\ShorthandsOff}{\@ifundefined{shorthandoff}{}{\shorthandoff{;:!?}}}
\newcommand{\ShorthandsOn}{\@ifundefined{shorthandon}{}{\shorthandon{;:!?}}}
\ShorthandsOff
\AtBeginDocument{\ShorthandsOn}
\AtBeginEnvironment{tikzpicture}{\ShorthandsOff}

\RequirePackage{geometry}
\RequirePackage{fancyhdr}
\RequirePackage{lastpage}
\RequirePackage{multicol}
\RequirePackage{array}
\RequirePackage{enumitem}
\RequirePackage{xcolor}
\RequirePackage{needspace}

\RequirePackage{amsmath,amssymb}
\RequirePackage{mathpazo}
\DeclareMathSizes{10.95}{10.95}{8}{5.5}
\begingroup\catcode`\_=\active \catcode`\^=\active
\gdef_#1{\sb{\scriptscriptstyle #1}}
\gdef^#1{\sp{\scriptscriptstyle\lycee@moinsepais #1}}
\endgroup
\newcommand\lycee@moins{\mathbin{%
  \setbox\z@\hbox{$\m@th\scriptscriptstyle\lycee@moinsorig$}%
  \dimen@=.55\wd\z@ \advance\dimen@-.5pt
  \hbox to.75\wd\z@{\hss$\m@th\scriptscriptstyle\vcenter{\hbox{%
    \tikz\draw[line width=.5pt,line cap=round](0,0)--(\the\dimen@,0);}}$\hss}}}
\begingroup\lccode`\~=`\- \lowercase{\endgroup\let~}\lycee@moins
\newcommand\lycee@moinsepais{\mathcode`\-="8000 }
\AtBeginDocument{\mathchardef\lycee@moinsorig=\mathcode`\-\relax}
\AtBeginDocument{\catcode`\_=12 \mathcode`\_="8000
  \catcode`\^=12 \mathcode`\^="8000 }
\newcommand\lycee@exposants{%
  \fontdimen13\textfont2=.48\fontdimen6\textfont2
  \fontdimen14\textfont2=.48\fontdimen6\textfont2
  \fontdimen15\textfont2=.48\fontdimen6\textfont2
  \fontdimen13\scriptfont2=.48\fontdimen6\scriptfont2
  \fontdimen14\scriptfont2=.48\fontdimen6\scriptfont2
  \fontdimen15\scriptfont2=.48\fontdimen6\scriptfont2}
\every@math@size\expandafter{\the\every@math@size\lycee@exposants}
\newlength{\retraitcalculs}
\setlength{\retraitcalculs}{2em}

\RequirePackage{tikz}
\usetikzlibrary{arrows.meta,calc,babel,shapes.misc}
\RequirePackage[most]{tcolorbox}
\tcbuselibrary{skins,breakable}

\definecolor{coulDef}{HTML}{1F4E79}
\definecolor{coulProp}{HTML}{7030A0}
\definecolor{coulRem}{HTML}{C55A11}
\definecolor{coulEx}{HTML}{2E7D32}
\definecolor{coulExo}{HTML}{B03A2E}
\definecolor{coulPart}{HTML}{1F4E79}
\definecolor{coulMeth}{HTML}{1B6E4F}
\definecolor{coulCourbe}{HTML}{0B5ED7}
\definecolor{coulCourbeB}{HTML}{C00000}
\definecolor{coulCourbeC}{HTML}{2E7D32}
\definecolor{coulGrille}{HTML}{8C8C8C}
\definecolor{coulRep}{HTML}{1B6E4F}
\definecolor{coulBar}{HTML}{2E7D32}

\newcommand{\R}{\mathbb{R}}
\newcommand{\Rs}{\mathbb{R}^{*}}
\renewcommand{\le}{\leqslant}
\renewcommand{\ge}{\geqslant}
\newcommand{\vect}[1]{\overrightarrow{#1}}
\newcommand{\norme}[1]{\left\lVert #1 \right\rVert}
\newcommand{\intervalle}[2]{\left[#1\,;#2\right]}
\RequirePackage{cancel}
\renewcommand{\CancelColor}{\color{coulExo}}
\newcommand{\fois}[1]{\times{\color{coulExo}#1}}
\newcommand{\barre}[1]{\cancel{#1}}

\tikzset{grille/.style={coulGrille,line width=0.4pt}}
\newcommand{\quadrillage}[4]{\draw[grille,step=1] (#1,#2) grid (#3,#4);}
\tikzset{
  croix/.style={cross out,draw,line width=0.6pt,inner sep=0pt,minimum size=4pt},
  croixdroite/.style={croix,rotate=45,line width=0.8pt},
}
\newcommand{\pt}[3]{\path (#1) node[croix]{} node[#2,font=\small,inner sep=2.5pt]{$#3$};}

\newcommand{\lycee@repere}[4]{%
  \draw[grille,step=1] (#1,#2) grid (#3,#4);
  \draw[-{Stealth[length=2mm]},thick] (#1,0)--({#3+0.4},0) node[below left=-1pt]{$x$};
  \draw[-{Stealth[length=2mm]},thick] (0,#2)--(0,{#4+0.4}) node[below left=-1pt]{$y$};
  \node[below left,font=\scriptsize,inner sep=1.5pt] at (0,0) {$0$};}
\newcommand{\lycee@gradx}[1]{\foreach \k/\l in {#1}{%
  \draw (\k,0.07)--(\k,-0.07) node[below,font=\scriptsize,inner sep=1.5pt]{$\l$};}}
\newcommand{\lycee@grady}[1]{\foreach \k/\l in {#1}{%
  \draw (0.07,\k)--(-0.07,\k) node[left,font=\scriptsize,inner sep=1.5pt]{$\l$};}}
\newcommand{\lycee@intff}[2]{\left[#1\,;#2\right]}
\newcommand{\lycee@intfo}[2]{\left[#1\,;#2\right[}
\newcommand{\lycee@intof}[2]{\left]#1\,;#2\right]}
\newcommand{\lycee@intoo}[2]{\left]#1\,;#2\right[}
\newcommand{\lycee@ens}[1]{\left\{#1\right\}}
\AtBeginDocument{%
  \providecommand{\repere}{\lycee@repere}%
  \providecommand{\gradx}{\lycee@gradx}%
  \providecommand{\grady}{\lycee@grady}%
  \providecommand{\Cf}{\mathcal{C}\sb{\scriptscriptstyle f}}%
  \providecommand{\Cg}{\mathcal{C}\sb{\scriptscriptstyle g}}%
  \providecommand{\intff}{\lycee@intff}%
  \providecommand{\intfo}{\lycee@intfo}%
  \providecommand{\intof}{\lycee@intof}%
  \providecommand{\intoo}{\lycee@intoo}%
  \providecommand{\ens}{\lycee@ens}}

\newlist{questions}{enumerate}{3}
\setlist[questions,1]{label=\arabic*\textdegree),leftmargin=*,itemsep=2pt,topsep=3pt}
\setlist[questions,2]{label=\alph*),leftmargin=*,itemsep=1pt,topsep=2pt}
\setlist[questions,3]{label=\roman*),leftmargin=*,itemsep=1pt,topsep=1pt}
\setlist[itemize]{label=\textbullet,leftmargin=*,itemsep=2pt,topsep=3pt}

\newcommand{\besoinplace}[1]{\par\penalty\@M\begingroup
  \dimen@=#1\relax\dimen@ii=\pagegoal\advance\dimen@ii-\pagetotal
  \ifdim\dimen@>\dimen@ii\ifdim\dimen@ii>\z@\vfil\fi\break\fi\endgroup}

\newcommand{\lycee@pied}{\thepage/\pageref{LastPage}}
\newcommand{\entete}[2]{%
  \pagestyle{fancy}\fancyhf{}%
  \fancyhead[L]{\footnotesize\color{coulPart}\textbf{#1}}%
  \fancyhead[R]{\footnotesize\color{coulPart}\textbf{#2}}%
  \fancyfoot[C]{\footnotesize\lycee@pied}%
  \renewcommand{\headrulewidth}{0.4pt}%
  \renewcommand{\headrule}{\hbox to\headwidth{%
    \color{coulPart!40}\leaders\hrule height \headrulewidth\hfill}}}

\RequirePackage{ccicons}
\newcommand{\lycee@licence}{{\scriptsize\color{gray}%
  \copyright\ Luc Destombe --- \ccbyncsa\ CC BY-NC-SA 4.0}}
\newif\iflycee@licence \lycee@licencetrue
\newcommand{\sanslicence}{\lycee@licencefalse}
\AtBeginDocument{\iflycee@licence\AddToHookNext{shipout/foreground}{%
  \put(0,0){\rlap{\kern\dimexpr1in+\hoffset+\oddsidemargin+\textwidth\relax
    \llap{\raisebox{-\dimexpr1in+\voffset+\topmargin+\headheight+\headsep
      +\textheight+\footskip\relax}{\lycee@licence}}}}}\fi}

\newcommand{\formulesserrees}{%
  \setlength{\abovedisplayskip}{3pt}\setlength{\belowdisplayskip}{3pt}%
  \setlength{\abovedisplayshortskip}{2pt}\setlength{\belowdisplayshortskip}{3pt}}

\setlength{\parindent}{0pt}

\RequirePackage[implicit=false,hidelinks,pdfencoding=auto,psdextra,bookmarksopen,
  bookmarksopenlevel=1,pdfcreator={},pdfproducer={}]{hyperref}
\RequirePackage{bookmark}
\hypersetup{pdfauthor={Luc Destombe},pdfkeywords={CC BY-NC-SA 4.0}}
\newcounter{lycee@signet}
\newcommand{\signet}[3][]{\stepcounter{lycee@signet}%
  \edef\lycee@dest{\if\relax\detokenize{#1}\relax
    signet.\arabic{lycee@signet}\else#1\fi}%
  \hypertarget{\lycee@dest}{}\bookmark[level=#2,dest=\lycee@dest]{#3}}
\pdfstringdefDisableCommands{%
  \def\R{R}\def\vect#1{#1}\def\Cf{Cf}\def\Cg{Cg}\def\fois#1{×#1}%
  \def\barre#1{#1}\def\intervalle#1#2{[#1;#2]}\def\intff#1#2{[#1;#2]}%
  \def\textsuperscript#1{#1}\def\dfrac#1#2{#1/#2}\def\frac#1#2{#1/#2}%
  \def\vec#1{vecteur #1}}
\RequirePackage[official]{eurosym}

\geometry{top=0.8cm,bottom=1.3cm,left=1cm,right=1cm,footskip=0.6cm}

\pagestyle{fancy}
\fancyhf{}
\renewcommand{\headrulewidth}{0pt}
\fancyfoot[C]{\footnotesize Page \thepage{} / \pageref{LastPage}}

\newcommand{\titrefiche}[2]{%
  \begin{center}
    \Large\textbf{#1}\\[2pt]
    \large\textbf{#2}\\[2pt]
  \end{center}
  \vspace{0.10cm}}

\setlength{\columnseprule}{0.4pt}
\setlength{\columnsep}{15pt}
\renewcommand{\columnseprulecolor}{\color{black!45}}
\setlength{\multicolsep}{2pt}

\newcounter{partiefiche}
\renewcommand{\thepartiefiche}{\Roman{partiefiche}}
\newif\ifapresbandeau
\newcommand{\lycee@bandeau}[1]{%
  \par\penalty-600
  \vspace{0.08cm}%
  \noindent\colorbox{coulPart}{%
    \parbox{\dimexpr\linewidth-2\fboxsep\relax}{%
      \color{white}\bfseries\raggedright\strut#1\strut}}%
  \nopagebreak\par\nopagebreak
  \vspace{0.06cm}\nopagebreak
  \apresbandeautrue}
\newcommand{\partiefiche}[1]{%
  \refstepcounter{partiefiche}%
  \lycee@bandeau{\signet{1}{\thepartiefiche{} --- #1}\thepartiefiche{} --- #1}}
\newcommand{\partiefichesansnum}[1]{\lycee@bandeau{\signet{1}{#1}#1}}

\newcounter{exo}
\newlength{\espaceexo}\setlength{\espaceexo}{7pt}
\newlength{\espacetitre}\setlength{\espacetitre}{2pt}
\newcommand{\exo}[1][\relax]{%
  \par
  \ifapresbandeau\apresbandeaufalse\else\penalty-150\addvspace{\espaceexo}\fi
  \refstepcounter{exo}%
  \noindent\signet[exercice-\arabic{exo}]{2}{Exercice \arabic{exo}}%
  \textbf{\color{coulExo}\underline{Exercice \theexo}}%
  \ifx\relax#1\relax\else\textbf{ : #1}\fi
  \par\nopagebreak\vspace{\espacetitre}\nopagebreak
  \ignorespaces}

\setlist[enumerate]{nosep,leftmargin=*,itemsep=2pt,topsep=2pt}
\setlist[itemize]{nosep,leftmargin=*,topsep=2pt}
\newcommand{\choix}[3]{\par\hspace*{1em}%
  \makebox[0.3\linewidth][l]{a) #1}\makebox[0.3\linewidth][l]{b) #2}c) #3}
\newcommand{\deux}[2]{\par\hspace*{0.6em}%
  \makebox[0.48\linewidth][l]{#1}#2}
\newcommand{\trois}[3]{\par\hspace*{0.6em}%
  \makebox[0.32\linewidth][l]{#1}\makebox[0.32\linewidth][l]{#2}#3}

\renewenvironment{center}{\par\addvspace{3pt}\centering}{\par\addvspace{3pt}}
\formulesserrees
\AtBeginDocument{\formulesserrees}
\clubpenalty=10000
\widowpenalty=10000
\displaywidowpenalty=10000
\brokenpenalty=10000
\setlength{\parskip}{0pt}

\RequirePackage{ragged2e}
\setlength{\RaggedRightRightskip}{0pt plus 1fil}
\AtBeginDocument{\RaggedRight\binoppenalty=10000 \relpenalty=10000 }
\g@addto@macro\@parboxrestore{\RaggedRight}
\makeatother
\usepackage{tkz-tab}
\geometry{bottom=1cm}

\setlength{\espaceexo}{4pt}
\setlist[enumerate]{nosep,leftmargin=*,itemsep=1pt,topsep=1pt}

\begin{document}

\titrefiche{Seconde --- Variations, signe et inéquations}{Fiche d'exercices du chapitre}

\begin{multicols}{2}\raggedcolumns

\partiefiche{Fonction croissante, fonction décroissante}

\exo[Lire les variations]
La courbe $\Cf$ d'une fonction $f$ est tracée ci-dessous.
\begin{center}
\begin{tikzpicture}[x=0.42cm,y=0.42cm]
  \repere{-4}{-3}{6}{4}
  \gradx{-3,-2,-1,1,2,3,4,5} \grady{-2,-1,1,2,3}
  \draw[coulCourbe,line width=1.1pt]
    (-3,-2) .. controls (-2,0.5) and (-1,3) .. (0,3)
            .. controls (1,3) and (2,-1) .. (3,-1)
            .. controls (3.67,-1) and (4.33,1) .. (5,2);
  \path[coulCourbe] (-3,-2) node[croix]{} (5,2) node[croix]{};
  \node[coulCourbe,font=\footnotesize] at (-2.6,1) {$\Cf$};
\end{tikzpicture}
\end{center}
\begin{enumerate}
  \item Sur quel intervalle $f$ est-elle définie ?
  \item Sur quels intervalles $f$ est-elle croissante ? décroissante ?
\end{enumerate}

\exo[Une randonnée]
À 8~h, Inès quitte le parking (altitude 1\,200~m) et monte jusqu'à un refuge, qu'elle
atteint à 11~h. Elle s'y repose jusqu'à 12~h, puis redescend au parking, où elle
arrive à 15~h. On note $f(t)$ son altitude à l'heure $t$, pour $t$ dans
$\intff{8}{15}$.
\begin{enumerate}
  \item Sur quels intervalles $f$ est-elle croissante ? constante ? décroissante ?
  \item Comparer $f(9)$ et $f(10)$, puis $f(13)$ et $f(14)$.
\end{enumerate}

\exo[Comparer des images]
\begin{enumerate}
  \item $g$ est une fonction définie sur $\R$. Comparer $g(2)$ et $g(5)$ lorsque :
        a)~$g$ est croissante ; b)~$g$ est décroissante ; c)~$g$ est constante.
  \item $f$ est croissante sur $\intff{-5}{5}$. Comparer $f(-4)$ et $f(-1)$.
\end{enumerate}

\exo[Conserver ou renverser l'ordre]
La fonction $f$ est croissante sur $\intff{0}{10}$ ; la fonction $g$ est décroissante
sur $\intff{0}{10}$.
\begin{enumerate}
  \item Ranger $f(1)$, $f(3)$, $f(4)$ et $f(8)$ dans l'ordre croissant.
  \item Ranger $g(1)$, $g(3)$, $g(4)$ et $g(8)$ dans l'ordre croissant.
  \item Recopier et compléter : « Une fonction croissante \dots{} l'ordre ; une
        fonction décroissante \dots{} l'ordre. »
  \item On sait que $f(a)<f(b)$. Peut-on avoir $a>b$ ? Justifier.
\end{enumerate}

\exo[Tracer une courbe]
Dans un repère, tracer une courbe passant par $A(0\,;1)$ et $B(6\,;3)$ et pouvant
représenter une fonction $f$ définie sur $\intff{0}{6}$ :
\begin{enumerate}
  \item croissante sur $\intff{0}{6}$ ;
  \item décroissante sur $\intff{0}{2}$, puis croissante sur $\intff{2}{6}$.
\end{enumerate}

\exo[Vrai ou faux]
Dire si chaque affirmation est vraie ou fausse, puis justifier (une courbe peut
servir de contre-exemple).
\begin{enumerate}
  \item Si $f$ est croissante sur $\intff{0}{10}$, alors elle est croissante sur
        $\intff{2}{5}$.
  \item Si $f(1)<f(4)$, alors $f$ est croissante sur $\intff{1}{4}$.
  \item Si $f$ n'est pas croissante sur $\intff{0}{5}$, alors elle est décroissante
        sur $\intff{0}{5}$.
\end{enumerate}

\partiefiche{Maximum, minimum}

\exo[Lire des extremums]
On reprend la fonction $f$ de l'exercice 1. Donner son maximum et son minimum, et où
ils sont atteints :
\begin{enumerate}
  \item sur $\intff{-3}{5}$ ;
  \item sur $\intff{0}{5}$ ;
  \item sur $\intff{3}{5}$.
\end{enumerate}

\exo[Tracer à partir d'une description]
Une fonction $f$ est définie sur $\intff{-2}{6}$. Elle est croissante sur
$\intff{-2}{1}$, décroissante sur $\intff{1}{4}$ et croissante sur $\intff{4}{6}$ ; son
maximum est $4$, atteint en $1$, et son minimum est $-3$, atteint en $4$. De plus,
$f(-2)=0$ et $f(6)=2$.
\begin{enumerate}
  \item Placer dans un repère les points connus de la courbe, puis tracer une courbe
        pouvant représenter $f$.
  \item Donner le maximum et le minimum de $f$ sur $\intff{4}{6}$.
\end{enumerate}

\exo[Les marées]
Un jour de grande marée, à Saint-Malo, la mer est haute à 2~h (hauteur d'eau
$12$~m), basse à 8~h ($2$~m), haute à 14~h ($13$~m) et basse à 20~h ($1$~m). Entre
deux de ces instants, la hauteur varie toujours dans le même sens. On note $f(t)$ la
hauteur d'eau (en m) à l'heure $t$, pour $t$ dans $\intff{2}{20}$.
\begin{enumerate}
  \item Donner les variations de $f$.
  \item Donner le maximum et le minimum de $f$ sur $\intff{2}{20}$, puis sur
        $\intff{2}{8}$, et à quelle heure ils sont atteints.
  \item Tracer une courbe pouvant représenter $f$ : en abscisse, l'heure de $0$ à
        $20$ (1 carreau pour 2~h) ; en ordonnée, la hauteur de $0$ à $14$~m
        (1 carreau pour 1~m). Placer d'abord les quatre points connus.
\end{enumerate}

\partiefiche{Tableau de variations}

\exo[De la courbe au tableau]
Dresser le tableau de variations de la fonction $f$ de l'exercice 1.

\exo[Du tableau à la courbe]
Une fonction $g$ a le tableau de variations ci-dessous.
\begin{center}
\begin{tikzpicture}[scale=0.85]
  \tkzTabInit[lgt=1.3,espcl=1.7,deltacl=0.6]{$x$/1,$g(x)$/2}{$-4$,$-1$,$2$,$5$}
  \tkzTabVar{+/$2$,-/$-3$,+/$3$,-/$-1$}
\end{tikzpicture}
\end{center}
\begin{enumerate}
  \item Sur quel intervalle $g$ est-elle définie ?
  \item Décrire ses variations par des phrases.
  \item Tracer, dans un repère, une courbe pouvant représenter $g$.
\end{enumerate}

\exo[Exploiter un tableau]
Une fonction $k$ a le tableau de variations ci-dessous.
\begin{center}
\begin{tikzpicture}[scale=0.85]
  \tkzTabInit[lgt=1.3,espcl=2.2,deltacl=0.6]{$x$/1,$k(x)$/2}{$-5$,$-1$,$4$}
  \tkzTabVar{-/$-2$,+/$4$,-/$0$}
\end{tikzpicture}
\end{center}
\begin{enumerate}
  \item Comparer $k(-4)$ et $k(-2)$, puis $k(1)$ et $k(3)$.
  \item Expliquer pourquoi on ne peut pas comparer $k(-3)$ et $k(2)$.
  \item Donner le maximum et le minimum de $k$.
\end{enumerate}

\exo[Extremums et tableau]
Une fonction $h$ a le tableau de variations ci-dessous.
\begin{center}
\begin{tikzpicture}[scale=0.85]
  \tkzTabInit[lgt=1.3,espcl=1.7,deltacl=0.6]{$x$/1,$h(x)$/2}{$-4$,$0$,$3$,$6$}
  \tkzTabVar{-/$2$,+/$5$,-/$-3$,+/$1$}
\end{tikzpicture}
\end{center}
\begin{enumerate}
  \item Donner le maximum et le minimum de $h$ sur $\intff{-4}{6}$, sur $\intff{-4}{0}$,
        puis sur $\intff{3}{6}$, et où ils sont atteints.
  \item Donner deux nombres $m$ et $M$ tels que $m\le h(x)\le M$ pour tout $x$ de
        $\intff{-4}{6}$.
\end{enumerate}

\exo[D'une description au tableau]
Une fonction $f$ est définie sur $\intff{-3}{6}$ ; elle est décroissante sur
$\intff{-3}{1}$, croissante sur $\intff{1}{4}$ et décroissante sur $\intff{4}{6}$. De
plus, $f(-3)=5$, $f(1)=-2$, $f(4)=3$ et $f(6)=0$.
\begin{enumerate}
  \item Dresser le tableau de variations de $f$.
  \item Donner son maximum et son minimum.
  \item Tracer, dans un repère, une courbe pouvant représenter $f$.
\end{enumerate}

\partiefiche{Fonctions affines et linéaires}

\exo[Reconnaître une fonction affine]
Pour chaque fonction, dire si elle est affine ; si oui, donner $a$ et $b$, et
préciser si elle est linéaire ou constante.
\deux{$f(x)=5x+2$}{$g(x)=6-2x$}
\deux{$h(x)=\frac{x}{3}$}{$k(x)=-4$}
\deux{$m(x)=x^{2}+3$}{$p(x)=\frac{x-1}{2}$}

\exo[Développer d'abord]
Développer et réduire, puis dire si la fonction est affine et donner ses variations.
\par\hspace*{0.6em}${f(x)=(x+3)^{2}-x^{2}}$
\par\hspace*{0.6em}${g(x)=(2x-1)(x+2)-2x^{2}}$
\par\hspace*{0.6em}${h(x)=(x+1)(x-1)-x(x+2)}$

\exo[Sens de variation]
\begin{enumerate}
  \item Donner le sens de variation de chaque fonction :
        \deux{$f(x)=-3x+1$}{$g(x)=0{,}5x-4$}
        \deux{$h(x)=7$}{$k(x)=2-\frac{x}{5}$}
  \item Sans calcul, comparer $f(-2)$ et $f(3)$, puis $g(-2)$ et $g(3)$.
\end{enumerate}

\exo[Tracer des droites]
Dans un même repère, tracer les droites représentant les fonctions
${f(x)=2x-1}$, ${g(x)=-x+3}$, ${h(x)=\frac{1}{2}x+1}$ et ${k(x)=-2}$.

\exo[Deux formules d'abonnement]
Une salle de sport propose deux formules mensuelles : A, $30$~€ par mois ; B,
$12$~€ par mois plus $3$~€ par séance. On note $x$ le nombre de séances dans le
mois, entre $0$ et $10$.
\begin{enumerate}
  \item Exprimer les prix $A(x)$ et $B(x)$ en fonction de $x$.
  \item Quelle est la nature de chaque fonction ? Donner leurs variations.
  \item Tracer leurs représentations dans un repère (1 carreau pour une séance
        en abscisse, 1 carreau pour 5~€ en ordonnée).
  \item Lire à partir de combien de séances la formule A est la moins chère.
  \item Le confirmer en calculant $B(6)$.
\end{enumerate}

\partiefiche{Fonction affine par lecture graphique}

\exo[Lire l'expression]
Les droites $(d_1)$ à $(d_4)$ représentent des fonctions affines $f_1$ à $f_4$.
\begin{center}
\begin{tikzpicture}[x=0.4cm,y=0.4cm]
  \repere{-4}{-4}{5}{5}
  \gradx{-3,-2,-1,1,2,3,4} \grady{-3,-2,-1,1,2,3,4}
  \draw[coulCourbe,line width=1.1pt] (-1,-4)--(3.5,5);
  \draw[coulCourbeB,line width=1.1pt] (-4,4)--(5,-0.5);
  \draw[coulCourbeC,line width=1.1pt] (-4,3)--(3,-4);
  \draw[black!70,line width=1.1pt] (-4,-2.667)--(5,3.333);
  \node[coulCourbe,font=\footnotesize,right] at (3.4,4.6) {$(d_1)$};
  \node[coulCourbeB,font=\footnotesize,above] at (-3.4,4.2) {$(d_2)$};
  \node[coulCourbeC,font=\footnotesize,left] at (-3.9,2.6) {$(d_3)$};
  \node[black!70,font=\footnotesize,right] at (5,3.3) {$(d_4)$};
\end{tikzpicture}
\end{center}
\begin{enumerate}
  \item Donner l'expression de $f_1$, $f_2$, $f_3$ et $f_4$.
  \item Laquelle est linéaire ?
\end{enumerate}

\exo[Associer fonctions et droites]
Associer chaque fonction à sa droite, en justifiant :
${f(x)=2x+1}$, ${g(x)=2x-3}$, ${h(x)=-x+1}$ et ${k(x)=-\frac{1}{3}x+1}$.
\begin{center}
\begin{tikzpicture}[x=0.4cm,y=0.4cm]
  \repere{-4}{-4}{4}{4}
  \gradx{-3,-2,-1,1,2,3} \grady{-3,-2,-1,1,2,3}
  \draw[coulCourbe,line width=1.1pt] (-3,4)--(4,-3);
  \draw[coulCourbeB,line width=1.1pt] (-2.5,-4)--(1.5,4);
  \draw[coulCourbeC,line width=1.1pt] (-4,2.333)--(4,-0.333);
  \draw[black!70,line width=1.1pt] (-0.5,-4)--(3.5,4);
  \node[coulCourbe,font=\footnotesize,left] at (-3,4) {$(D_1)$};
  \node[coulCourbeB,font=\footnotesize,right] at (1.5,4) {$(D_2)$};
  \node[coulCourbeC,font=\footnotesize,above] at (-3.4,2.3) {$(D_3)$};
  \node[black!70,font=\footnotesize,right] at (3.5,4) {$(D_4)$};
\end{tikzpicture}
\end{center}

\exo[Location de vélo]
La droite ci-dessous donne le prix $p(x)$, en euros, de la location d'un vélo
électrique pendant $x$ heures.
\begin{center}
\begin{tikzpicture}[x=0.55cm,y=0.13cm]
  \draw[grille,xstep=1,ystep=2] (0,0) grid (7,24);
  \draw[-{Stealth[length=2mm]},thick] (0,0)--(7.5,0) node[right,font=\footnotesize]{durée (h)};
  \draw[-{Stealth[length=2mm]},thick] (0,0)--(0,25.5) node[right,font=\footnotesize]{prix (€)};
  \node[below left,font=\scriptsize,inner sep=1.5pt] at (0,0) {$0$};
  \foreach \k in {1,...,7} \draw (\k,0.4)--(\k,-0.4) node[below,font=\scriptsize,inner sep=1.5pt]{$\k$};
  \foreach \k in {4,8,...,24} \draw (0.1,\k)--(-0.1,\k) node[left,font=\scriptsize,inner sep=1.5pt]{$\k$};
  \draw[coulCourbe,line width=1.1pt] (0,4)--(6.667,24);
\end{tikzpicture}
\end{center}
\begin{enumerate}
  \item Lire le prix pour 2~h de location.
  \item Lire l'ordonnée à l'origine $b$. Que représente-t-elle ?
  \item Lire le coefficient directeur $a$. Que représente-t-il ?
  \item En déduire $p(x)$, puis calculer le prix pour 8~h.
\end{enumerate}

\partiefiche{Fonction affine passant par deux points}

\exo[Deux images connues]
Trouver la fonction affine qui vérifie :
\deux{$f(1)=4$ et $f(3)=10$}{$g(-2)=7$ et $g(2)=-1$}
\deux{$h(0)=5$ et $h(4)=3$}{$k(-3)=-4$ et $k(6)=2$}

\exo[Droite passant par deux points]
\begin{enumerate}
  \item Trouver la fonction affine $f$ dont la droite passe par $A(2\,;1)$ et
        $B(6\,;3)$. Que remarque-t-on ?
  \item Le point $E(10\,;5)$ est-il sur la droite $(AB)$ ?
  \item Trouver la fonction affine $g$ dont la droite passe par $C(-1\,;5)$ et
        $D(2\,;-4)$.
  \item Le point $F(3\,;-6)$ est-il sur la droite $(CD)$ ?
\end{enumerate}

\exo[Température]
Un matin de printemps, il fait $6$~°C à 8~h et $18$~°C à 14~h. On admet que la
température $T(h)$ est une fonction affine de l'heure $h$ sur $\intff{8}{14}$.
\begin{enumerate}
  \item Donner $T(8)$ et $T(14)$, puis l'expression de $T(h)$.
  \item Calculer la température à 11~h.
  \item À quelle heure fait-il $14$~°C ?
\end{enumerate}

\exo[Facture d'électricité]
Une facture d'électricité est une fonction affine $f$ de la consommation $x$ (en
kWh) : pour $300$~kWh, on paie $72$~€ ; pour $500$~kWh, on paie $102$~€.
\begin{enumerate}
  \item Trouver $f(x)$.
  \item Que représentent $a$ et $b$ ?
  \item Combien paie-t-on pour $800$~kWh ?
\end{enumerate}

\partiefiche{Signe d'une fonction par lecture graphique}

\exo[Signe d'une fonction]
La courbe $\Cf$ d'une fonction $f$ définie sur $\intff{-4}{5}$ est tracée ci-dessous.
\begin{center}
\begin{tikzpicture}[x=0.42cm,y=0.42cm]
  \repere{-5}{-3}{6}{3}
  \gradx{-4,-3,-2,-1,1,2,3,4,5} \grady{-2,-1,1,2}
  \draw[coulCourbe,line width=1.1pt]
    (-4,2) .. controls (-3.67,1.33) and (-3.33,0.44) .. (-3,0)
           .. controls (-2.33,-0.89) and (-1.67,-2) .. (-1,-2)
           .. controls (-0.33,-2) and (0.33,-0.89) .. (1,0)
           .. controls (1.33,0.44) and (1.67,2) .. (2,2)
           .. controls (2.67,2) and (3.33,0.89) .. (4,0)
           .. controls (4.33,-0.44) and (4.67,-1.33) .. (5,-2);
  \path[coulCourbe] (-4,2) node[croix]{} (5,-2) node[croix]{};
  \node[coulCourbe,font=\footnotesize] at (-4.4,1) {$\Cf$};
\end{tikzpicture}
\end{center}
\begin{enumerate}
  \item Résoudre graphiquement $f(x)=0$.
  \item Dresser le tableau de signes de $f$.
  \item Résoudre graphiquement $f(x)>0$.
  \item Dresser le tableau de variations de $f$.
\end{enumerate}

\exo[Signe de fonctions affines]
Dresser le tableau de signes des fonctions $f_1$, $f_2$, $f_3$ et $f_4$ de
l'exercice 20.

\exo[Des variations au signe]
Une fonction $f$, définie sur $\intff{-3}{5}$, a le tableau de variations ci-dessous ;
de plus, $f(-1)=0$ et $f(3)=0$. Dresser son tableau de signes.
\begin{center}
\begin{tikzpicture}[scale=0.85]
  \tkzTabInit[lgt=1.3,espcl=2.2,deltacl=0.6]{$x$/1,$f(x)$/2}{$-3$,$1$,$5$}
  \tkzTabVar{-/$-2$,+/$4$,-/$-1$}
\end{tikzpicture}
\end{center}

\partiefiche{Propriétés des inégalités, encadrer une quantité}

\exo[Règles des inégalités]
On sait que $a<b$. Compléter par $<$ ou $>$.
\deux{a)~$a+4\ \dots\ b+4$}{b)~$a-7\ \dots\ b-7$}
\deux{c)~$3a\ \dots\ 3b$}{d)~$-5a\ \dots\ -5b$}
\deux{e)~$\dfrac{a}{2}\ \dots\ \dfrac{b}{2}$}{f)~$-a\ \dots\ -b$}

\exo[Encadrer]
On sait que $1\le x\le 4$. Encadrer chaque expression.
\deux{$A=2x+3$}{$B=-3x+1$}
\deux{$C=5-2x$}{$D=\frac{1}{2}x-1$}

\exo[Mesures imprécises]
\begin{enumerate}
  \item Un potager rectangulaire mesure $6$~m de large ; sa longueur $x$, en
        mètres, vérifie $10\le x\le 12$. Encadrer son périmètre $P=2x+12$, puis son
        aire $\mathcal{A}=6x$.
  \item Une température $C$ en degrés Celsius vaut $F=1{,}8C+32$ en degrés
        Fahrenheit. Au printemps, $15\le C\le 25$. Encadrer $F$.
\end{enumerate}

\partiefiche{Résoudre une inéquation du premier degré}

\exo[Est-ce une solution ?]
Le nombre $3$ est-il solution de chaque inéquation ? Justifier.
\deux{a)~$2x-1>4$}{b)~$-x+5\ge 2$}
\deux{c)~$4x-9<x$}{d)~$5-2x\le -1$}

\exo[Résoudre]
Résoudre chaque inéquation ; donner $S$ sous forme d'intervalle.
\deux{a)~$x+6>2$}{b)~$-3x<15$}
\deux{c)~$2x-5\ge 7$}{d)~$9-4x>1$}
\deux{e)~$3x+2\le 5x-6$}{f)~$-x+4\ge 3x-8$}

\exo[Parenthèses et fractions]
Résoudre chaque inéquation.
\deux{a)~$2(x-3)<x+1$}{b)~$3(2-x)\ge 2(x+8)$}
\deux{c)~$\dfrac{x}{2}+1>4$}{d)~$\dfrac{2x-1}{3}\le x$}

\exo[Corriger les erreurs]
Chaque résolution ci-dessous est fausse. Trouver l'erreur, puis corriger.
\begin{enumerate}
  \item $-2x+3>9$, donc $-2x>6$, donc $x>-3$.
  \item $5x\le -10$, donc $x\ge -2$.
\end{enumerate}

\partiefiche{Signe d'une fonction affine par le calcul}

\exo[Tableaux de signes]
Dresser le tableau de signes de chaque expression.
\deux{$A(x)=2x-8$}{$B(x)=-3x+9$}
\deux{$C(x)=5x+15$}{$D(x)=-\frac{1}{2}x+4$}
\deux{$E(x)=7x$}{$F(x)=-6x+4$}

\exo[Variations, droite et signe]
On considère ${f(x)=3x-6}$ et ${g(x)=8-2x}$.
\begin{enumerate}
  \item Donner le sens de variation de $f$ et de $g$.
  \item Tracer leurs droites dans un même repère.
  \item Dresser le tableau de signes de $f$ et de $g$ par le calcul, puis le
        vérifier sur le graphique.
\end{enumerate}

\exo[Associer]
Associer chaque fonction ${f(x)=4x-12}$ et ${g(x)=6-2x}$ à son tableau de signes,
en justifiant.
\begin{center}
\begin{tikzpicture}[scale=0.8]
  \tkzTabInit[lgt=2,espcl=1.6]{$x$/1,Tableau 1/1}{$-\infty$,$3$,$+\infty$}
  \tkzTabLine{,+,z,-,}
\end{tikzpicture}\quad
\begin{tikzpicture}[scale=0.8]
  \tkzTabInit[lgt=2,espcl=1.6]{$x$/1,Tableau 2/1}{$-\infty$,$3$,$+\infty$}
  \tkzTabLine{,-,z,+,}
\end{tikzpicture}
\end{center}

\exo[Bénéfice]
Une artisane fabrique des bougies. Pour $x$ bougies vendues dans le mois (de $0$ à
$100$), son bénéfice, en euros, est ${B(x)=12x-480}$.
\begin{enumerate}
  \item Calculer $B(0)$. Que représente ce nombre ?
  \item Dresser le tableau de signes de $B(x)$ sur $\intff{0}{100}$.
  \item À partir de combien de bougies vendues gagne-t-elle de l'argent ?
\end{enumerate}

\partiefiche{Inéquations produits}

\exo[Signe d'un produit]
Dresser le tableau de signes de chaque produit.
\par\hspace*{0.6em}$A(x)=(x-2)(3x+6)$
\par\hspace*{0.6em}$B(x)=(4x-12)(-x+1)$
\par\hspace*{0.6em}$C(x)=(5-x)(2x+8)$

\exo[Résoudre]
Résoudre chaque inéquation.
\deux{$(x-4)(x+1)<0$}{$(2x+6)(x-5)\ge 0$}
\deux{$(3-x)(x+2)>0$}{$(-2x+4)(3x-9)\le 0$}

\exo[Trouver l'erreur]
Pour étudier le signe de $(x+1)(2-x)$, Lina a dressé le tableau ci-dessous. Trouver
son erreur, puis dresser le bon tableau.
\begin{center}
\begin{tikzpicture}[scale=0.8]
  \tkzTabInit[lgt=2.6,espcl=1.6]{$x$/1,$x+1$/1,$2-x$/1,$(x+1)(2-x)$/1}%
    {$-\infty$,$2$,$-1$,$+\infty$}
  \tkzTabLine{,-,t,-,z,+,}
  \tkzTabLine{,+,z,-,t,-,}
  \tkzTabLine{,-,z,+,z,-,}
\end{tikzpicture}
\end{center}

\partiefiche{Inéquations quotients}

\exo[Signe d'un quotient]
Pour chaque quotient, donner la valeur interdite, puis dresser le tableau de signes.
\trois{$A(x)=\dfrac{x+3}{x-2}$}{$B(x)=\dfrac{6-2x}{x+4}$}{$C(x)=\dfrac{3x}{5-x}$}

\exo[Résoudre]
Résoudre chaque inéquation.
\deux{$\dfrac{x-1}{x+3}\ge 0$}{$\dfrac{2x+8}{x-6}<0$}\smallskip
\deux{$\dfrac{4-x}{2x-2}\le 0$}{$\dfrac{-3x-6}{x+5}>0$}

\exo[Produit, puis quotient]
\begin{enumerate}
  \item Résoudre $(3x+6)(-x+4)>0$.
  \item Pour quelle valeur de $x$ le quotient $\dfrac{3x+6}{-x+4}$ n'existe-t-il pas ?
  \item Résoudre $\dfrac{3x+6}{-x+4}\le 0$.
\end{enumerate}

\partiefiche{Inéquations avancées}

\exo[Factoriser d'abord]
Résoudre chaque inéquation.
\trois{$x^{2}\ge 3x$}{$2x^{2}<8x$}{$x^{2}\le -5x$}

\exo[Avec une identité remarquable]
Résoudre chaque inéquation.
\trois{$x^{2}<25$}{$(x-3)^{2}\ge 4$}{$(2x+1)^{2}\le 9$}

\exo[Comparer un quotient à un nombre]
Résoudre chaque inéquation.
\deux{$\dfrac{x+5}{x-1}\ge 2$}{$\dfrac{3}{x+2}<1$}

\columnbreak
\exo[Coût moyen]
Une imprimante 3D fabrique des figurines : $x$ figurines coûtent $2x+60$~€ (matière
et réglages). Le coût moyen d'une figurine est ${M(x)=\dfrac{2x+60}{x}}$, pour
$x>0$.
\begin{enumerate}
  \item Calculer $M(10)$ et $M(30)$.
  \item Montrer que $M(x)\le 5$ revient à $\dfrac{60-3x}{x}\le 0$.
  \item À partir de combien de figurines le coût moyen est-il d'au plus $5$~€ ?
\end{enumerate}

\partiefiche{Problèmes d'inéquations}

\exo[Moyenne trimestrielle]
Une élève a obtenu $9$, $12$ et $10$ sur $20$ aux trois premiers devoirs du trimestre.
On note $n$ sa note au quatrième devoir.
\begin{enumerate}
  \item Exprimer la moyenne $M(n)$ de ses quatre notes.
  \item Quelles notes doit-elle obtenir pour avoir au moins $12$ de moyenne ?
  \item Peut-elle encore obtenir au moins $15$ de moyenne ?
\end{enumerate}

\exo[Rectangle ou carré]
Un rectangle a pour longueur $x+5$ et pour largeur $3$ ; un carré a pour côté $x$
(longueurs en cm, $x>0$).
\begin{enumerate}
  \item Exprimer le périmètre $P(x)$ du rectangle et le périmètre $Q(x)$ du carré.
  \item Pour quelles valeurs de $x$ le périmètre du carré est-il plus grand que
        celui du rectangle ?
  \item Vérifier avec $x=10$.
\end{enumerate}

\exo[Trottinettes en libre-service]
Deux opérateurs louent des trottinettes. Opérateur A : $1$~€ pour débloquer, puis
$0{,}20$~€ par minute ; opérateur B : $0{,}30$~€ par minute, sans déblocage. On note
$x$ la durée du trajet en minutes, entre $0$ et $30$.
\begin{enumerate}
  \item Exprimer les prix $A(x)$ et $B(x)$.
  \item Calculer le prix d'un trajet de $5$ minutes, puis de $20$ minutes, avec
        chaque opérateur.
  \item Résoudre l'inéquation $A(x)<B(x)$ sur $\intff{0}{30}$.
  \item Quel opérateur choisir, selon la durée du trajet ?
\end{enumerate}

\partiefichesansnum{Problème de synthèse}

\exo[Étude complète d'une fonction]
On considère la fonction $f$ définie sur $\intff{0}{6}$ par ${f(x)=(x-1)(x-5)}$.
\begin{enumerate}
  \item Dresser le tableau de valeurs de $f$ de $x=0$ à $x=6$, avec un pas de $1$.
  \item Tracer la courbe de $f$ dans un repère orthonormé (unité : 1 carreau).
  \item Donner les variations de $f$, puis dresser son tableau de variations.
  \item Donner le maximum et le minimum de $f$ sur $\intff{0}{6}$.
  \item Dresser, par le calcul, le tableau de signes de $f(x)$ sur $\intff{0}{6}$.
  \item Résoudre $f(x)<0$, puis le vérifier sur la courbe.
\end{enumerate}

\end{multicols}

\end{document}
